% Auto-generated by scripts/exp_fold_gauge_anomaly_clt_variance.py
\leanverified{paper\_fold\_gauge\_anomaly\_clt}
\begin{theorem}[规范差的中心极限定理与方差闭式]\label{thm:fold-gauge-anomaly-clt}
在均匀基线 $\omega\sim\mathrm{Unif}(\Omega_{m+1})$ 下，令 $G_m(\omega)$ 为规范差（定义 \ref{def:fold-gauge-anomaly}）。
则存在常数 $g_\ast,\sigma_G^2>0$ 使得：
\[
\frac{G_m-g_\ast m}{\sqrt m}\Rightarrow \mathcal N(0,\sigma_G^2),
\qquad
\mathrm{Var}(G_m)=\sigma_G^2\,m+o(m).
\]
并且这两个常数具有完全闭式：
\[
\boxed{\ g_\ast=\frac49,\qquad \sigma_G^2=\frac{118}{243}\ }.
\]
\end{theorem}

\begin{proof}[证明要点（可审计）]
用 Berstel 的 $4$ 状态不歧义 Zeckendorf 归一化换能器（\cite{ITA_2001__35_6_491_0}，Figure~1）描述逐位规范化。在 IID Bernoulli$(1/2)$ 输入下，唯一接受路径诱导一个原始有限状态 Markov 链；而规范差的逐位指示 $g_t=\mathbf 1\{X_t\neq Y_t\}$ 是该链上的可加观测量。因此由有限状态 Markov 链的 CLT 得到上述极限定理。
常数可由压力函数的导数闭式读出：对每条边赋权$w_\theta=e^{\theta\,\mathbf 1\{X\neq Y\}}\cdot\mathbb P(X)$，得到加权邻接矩阵 $A_\theta$；则 $P(\theta)=\log\rho(A_\theta)$ 的一二阶导数给出 $g_\ast=P'(0)$ 与 $\sigma_G^2=P''(0)$。脚本 \texttt{scripts/exp\_fold\_gauge\_anomaly\_clt\_variance.py} 以严格有理线性代数实现该计算并输出上述闭式。
\end{proof}
